\documentclass[11pt]{article}

\usepackage{amsmath,amssymb}
\usepackage{xurl}
\usepackage[hidelinks]{hyperref}
\setlength{\emergencystretch}{2em}

% Shared commands for notes in docs/paper-gaps/.

\DeclareUnicodeCharacter{2080}{\ifmmode{}_{0}\else\textsubscript{0}\fi}
\DeclareUnicodeCharacter{2081}{\ifmmode{}_{1}\else\textsubscript{1}\fi}
\DeclareUnicodeCharacter{2082}{\ifmmode{}_{2}\else\textsubscript{2}\fi}
\DeclareUnicodeCharacter{2099}{\ifmmode{}_{n}\else\textsubscript{n}\fi}

\newcommand{\Lean}{\textsc{Lean}}
\ExplSyntaxOn
\NewDocumentCommand{\leanid}{m}{
  \ifmmode
    \textsf{\detokenize{#1}}
  \else
    \group_begin:
    \urlstyle{sf}
    \tl_set:Nn \l_tmpa_tl {#1}
    \tl_replace_all:Nnn \l_tmpa_tl {\_} {_}
    \exp_args:NV \path \l_tmpa_tl
    \group_end:
  \fi
}
\ExplSyntaxOff
\providecommand{\tr}{\operatorname{tr}}
\newcommand{\tnleanIssueBase}{https://github.com/LionSR/TNLean/issues/}
\newcommand{\tnleanPrBase}{https://github.com/LionSR/TNLean/pull/}
\newcommand{\tnleanDocBase}{https://sirui-lu.com/TNLean/docs/}
\newcommand{\ghissue}[1]{\href{\tnleanIssueBase#1}{GitHub issue \##1}}
\newcommand{\ghpr}[1]{\href{\tnleanPrBase#1}{PR \##1}}
\newcommand{\doclink}[1]{\href{\tnleanDocBase#1.html}{\path{#1}}}

% Dirac notation and the identity operator, as in blueprint/src/macros/common.tex.
% \providecommand keeps note-local definitions authoritative.
\usepackage{dsfont}
\providecommand{\ket}[1]{|#1\rangle}
\providecommand{\bra}[1]{\langle #1|}
\providecommand{\braket}[2]{\langle #1 | #2 \rangle}
\providecommand{\ketbra}[2]{|#1\rangle\!\langle #2|}
\providecommand{\Id}{\mathds{1}}
\providecommand{\one}{\mathds{1}}
\providecommand{\C}{\mathbb{C}}
\providecommand{\MN}[1]{M_{#1}(\C)}
\providecommand{\MD}{M_D(\C)}

% External referencing for paper sources.  The build helper writes sanitized
% auxiliary files under build/paper-aux/ and excludes duplicated source labels.
\usepackage{xr}

% Import a generated paper auxiliary file with a caller-supplied label prefix.
% The candidate paths support builds from docs/paper-gaps/, the repository root,
% and build/paper-aux/ itself.
\newcommand{\importpaperaux}[2]{%
  \IfFileExists{../../build/paper-aux/#2.aux}{%
    \externaldocument[#1]{../../build/paper-aux/#2}%
  }{%
    \IfFileExists{build/paper-aux/#2.aux}{%
      \externaldocument[#1]{build/paper-aux/#2}%
    }{%
      \IfFileExists{#2.aux}{%
        \externaldocument[#1]{#2}%
      }{}%
    }%
  }%
}

% General external reference macros
\newcommand{\paperref}[2]{\ref{#1-#2}}
\newcommand{\papereqref}[2]{(\ref{#1-#2})}

% CPSV16 source (1606.00608 / MPDO-22-12-17-2.tex) external document and shortcuts
\importpaperaux{cpsv16-}{1606.00608/MPDO-22-12-17-2-tnlean}
\newcommand{\cpsvref}[1]{\paperref{cpsv16}{#1}}
\newcommand{\cpsveqref}[1]{\papereqref{cpsv16}{#1}}

% Verdict marker: \gapnote{<kind>}{<status>}. Kind is the policy
% classification (clarification, local-correction, scope-restriction,
% unfaithful, false-source, open-gap); status is open, wip (elimination
% underway), resolved, or historical. Machine-read by the site index;
% prints a verdict line.
\newcommand{\gapnote}[2]{\par\noindent\textbf{Verdict:} #1 (#2).\par}


\title{Scope of the Formal Block-Separation Lemmas Relative to PGVWC07 Uniqueness}
\author{}
\date{2026-05-12, updated 2026-09-05}

\begin{document}
\maketitle
\gapnote{scope-restriction}{resolved}

\section{The source theorem}

Let $B=(B_i)_i$ and $C=(C_i)_i$ be two translation-invariant canonical
$D$-dimensional MPS representations of the same state on a finite ring. This
notation follows the uniqueness theorem of P\'erez-Garc\'ia, Verstraete, Wolf,
and Cirac
\cite[Theorem~\texttt{thm-uniq}, lines~1002--1015]{PerezGarcia2007MPSReps}.

For a translation-invariant tensor $A=(A_i)_i$ and a length $L$, define
\begin{align}
    \Gamma_L(X)
     & :=
    \sum_{i_1,\ldots,i_L}
    \tr\!\left(XA_{i_1}\cdots A_{i_L}\right)
    \ket{i_1\cdots i_L}.
    \label{eq:pgvwc07_ti_uniqueness_scope_gamma_map}
\end{align}
Condition C1 in the source requires a finite length $L_0$ such that
\begin{align}
    \Gamma_{L_0}
     & \text{ is injective}.
    \label{eq:pgvwc07_ti_uniqueness_scope_finite_injectivity}
\end{align}
The source introduces this condition in lines~888--908 of
Ref.~\cite{PerezGarcia2007MPSReps}.

The uniqueness theorem assumes
(\ref{eq:pgvwc07_ti_uniqueness_scope_finite_injectivity}), uniqueness of the
open-boundary-condition (OBC) canonical representation of the MPS state, and
the finite-size bound
\begin{align}
    N
     & > 2L_0+D^4.
    \label{eq:pgvwc07_ti_uniqueness_scope_finite_size_bound}
\end{align}
As printed, it concludes that one unitary matrix $U$ intertwines the two
representations exactly:
\begin{align}
    B_i
     & =UC_iU^\dagger
    \qquad\text{for every }i.
    \label{eq:pgvwc07_ti_uniqueness_scope_printed_unitary_intertwiner}
\end{align}
When the Schmidt weights are nondegenerate, the source gives the usual
permutation-and-phase interpretation
\cite[Theorem~\texttt{thm-uniq}, lines~1002--1015]{PerezGarcia2007MPSReps}.

The exact conclusion
(\ref{eq:pgvwc07_ti_uniqueness_scope_printed_unitary_intertwiner}) overlooks a
finite-ring phase. If the tensors represent the same nonzero length-$N$ state,
the corrected conclusion has a unitary $U$ and a scalar $\xi\in\C$ satisfying
\begin{align}
    |\xi|
     & =1,
    \label{eq:pgvwc07_ti_uniqueness_scope_unit_phase} \\
    \xi^N
     & =1,
    \label{eq:pgvwc07_ti_uniqueness_scope_root_phase} \\
    B_i
     & =\xi UC_iU^\dagger
    \qquad\text{for every }i.
    \label{eq:pgvwc07_ti_uniqueness_scope_corrected_unitary_intertwiner}
\end{align}
The counterexample and corrected reconstruction statement are recorded in
\path{docs/paper-gaps/pgvwc07_quadratic_reconstruction_phase.tex}.

\section{Formal boundary}

The checked Chapter~9 block-separation lemmas do not formalize
Theorem~\texttt{thm-uniq} of Ref.~\cite{PerezGarcia2007MPSReps}. They compare
block families that already share a block index set and block-dimension
function. Their substantive hypotheses are blockwise injectivity of the first
family and equality of the corresponding MPV families; the global and
canonical-form variants retain the same block correspondence and use a common,
otherwise arbitrary, weight function. They assume neither blockwise
normalization nor strict ordering of weight moduli.\footnote{The formal declarations are
    \leanid{MPSTensor.fundamentalTheorem_multiBlock_blocks},
    \leanid{MPSTensor.fundamentalTheorem_multiBlock_global},
    \leanid{MPSTensor.fundamentalTheorem_canonicalForm_sameStructure}, and
    \leanid{MPSTensor.fundamentalTheorem_multiBlock_explicit}.}

These same-structure assumptions are useful in local MPS arguments, but they
are not the source hypotheses. In particular, the source theorem does not
assume that the representations have a common prepared block decomposition.
Its proof instead uses
uniqueness of the OBC canonical representation and the Horn--Johnson
matrix-lemma route described after Theorem~\texttt{thm-uniq} in
Ref.~\cite{PerezGarcia2007MPSReps}.

The distinction is theorem-level. The same-structure results establish
conditional consequences under stronger assumptions and must not be cited as
formalizations of PGVWC07 Theorem~\texttt{thm-uniq}; the corrected theorem is
formalized separately, as recorded below.

\section{Required formalization}

A mathematically corrected theorem must quantify over two
translation-invariant canonical $D$-MPS representations of the same nonzero
finite-ring state. It must assume
(\ref{eq:pgvwc07_ti_uniqueness_scope_finite_injectivity}), uniqueness of the
OBC canonical representation, and
(\ref{eq:pgvwc07_ti_uniqueness_scope_finite_size_bound}). Its conclusion must
provide a unitary $U$ and an $N$th-root phase $\xi$ satisfying
(\ref{eq:pgvwc07_ti_uniqueness_scope_unit_phase})--
(\ref{eq:pgvwc07_ti_uniqueness_scope_corrected_unitary_intertwiner}). Exact
unitary conjugacy holds only after the explicit rephasing
$\widehat C_i:=\xi C_i$, which preserves the length-$N$ state because
$\xi^N=1$.

The current same-structure lemmas may support such a proof only after the
source argument has produced a common block correspondence and established,
block by block, equality of the corresponding matrix product vector families.
They do not themselves derive that correspondence from finite-ring state
equality and therefore cannot replace the source theorem.

\section{Formal route}

The corrected theorem is now formalized. The declaration
\leanid{MPSTensor.pgvwc07_uniqueness_of_canonical_form} takes two
translation-invariant canonical representations $B$ and $C$ in the block
form of Theorem~4 of Ref.~\cite{PerezGarcia2007MPSReps}, that is, weighted
direct sums $\bigoplus_j\lambda_jA^j$ with $1\ge\lambda_j>0$ whose blocks are
unital, carry a positive-definite diagonal dual fixed point, and have scalar
fixed points of their transfer maps (the predicate
\leanid{MPSTensor.IsPGVWC07CanonicalForm}), condition
(\ref{eq:pgvwc07_ti_uniqueness_scope_finite_injectivity}) for $B$, the bound
(\ref{eq:pgvwc07_ti_uniqueness_scope_finite_size_bound}), and equality of the
two nonzero length-$N$ states, and concludes
(\ref{eq:pgvwc07_ti_uniqueness_scope_unit_phase})--%
(\ref{eq:pgvwc07_ti_uniqueness_scope_corrected_unitary_intertwiner}). The
companion \leanid{MPSTensor.pgvwc07_uniqueness_of_canonical_form_rephased}
states exact unitary conjugacy after the rephasing $\widehat C_i=\xi C_i$
together with the invariance of the length-$N$ state, and
\leanid{MPSTensor.pgvwc07_uniqueness_of_canonical_form_unitGaugePhaseEquiv}
packages the pair $(U,\xi)$ as a unit-modulus gauge-phase equivalence. The
single-block cases, with unit weights and with the block weights
$0<\lambda_B,\lambda_C\le1$ of Theorem~4, are
\leanid{MPSTensor.pgvwc07_uniqueness_of_unital} and
\leanid{MPSTensor.pgvwc07_uniqueness_of_weighted}; the latter concludes
$\lambda_B=\lambda_C$ as well.

Condition (\ref{eq:pgvwc07_ti_uniqueness_scope_finite_injectivity}) leaves
$B$ with a single block, the first assertion of the proposition on condition
C1 (lines~911--919;
\leanid{MPSTensor.PGVWC07CanonicalFormData.not_isNBlkInjective_of_ne}): a
matrix unit supported on an off-diagonal block is annihilated by every word.
The intertwiner $R$ obtained below is invertible by the same condition
(lines~1105--1108), so $C_i=x^{-1}R^{-1}B_iR$ is conjugate to a multiple of
$B_i$ and satisfies
(\ref{eq:pgvwc07_ti_uniqueness_scope_finite_injectivity}) as well; hence $C$
also has a single block, and the weighted single-block theorem applies. The
printed proof uses $\sum_iC_iC_i^\dagger=\Id$ at line~1104 and thereby treats
only a single unit-weight block for $C$; the block form of $C$ is handled by
this reduction.

The printed proof (lines~1063--1108) obtains invertible intertwiners on the
doubled bond space by rewriting both representations as open-boundary
representations and comparing them, through the freedom theorem
(Theorem~2, lines~466--486) and the uniqueness of the open-boundary
canonical form, with that canonical form. The open-boundary canonical form
is not formalized. The formal route obtains the same intertwiners directly.
For a cut after $k$ sites with $m=N-k$ sites to its right, the cut vectors
of the proof of the proposition on condition C1 (lines~921--931),
\begin{align}
    \lvert\Phi_{\alpha,\beta}\rangle
     & =\sum_{i_1,\ldots,i_k}\langle\alpha|A_{i_1}\cdots A_{i_k}|\beta\rangle
    \,\lvert i_1\cdots i_k\rangle,
    \label{eq:pgvwc07_ti_uniqueness_scope_left_cut_vector}                    \\
    \lvert\Psi_{\alpha,\beta}\rangle
     & =\sum_{j_1,\ldots,j_m}\langle\beta|A_{j_1}\cdots A_{j_m}|\alpha\rangle
    \,\lvert j_1\cdots j_m\rangle,
    \label{eq:pgvwc07_ti_uniqueness_scope_right_cut_vector}
\end{align}
resolve the ring state as
$\sum_{\alpha,\beta}\lvert\Phi_{\alpha,\beta}\rangle\lvert\Psi_{\alpha,\beta}\rangle$
(lines~921--924), and condition
(\ref{eq:pgvwc07_ti_uniqueness_scope_finite_injectivity}) makes both families
linearly independent whenever $k,m\geq L_0$ (lines~936--949; formal
declarations \leanid{MPSTensor.linearIndependent_leftCutVector} and
\leanid{MPSTensor.linearIndependent_rightCutVector}). Two decompositions of
the same state with linearly independent factors on the $B$ side are related
by an invertible matrix $W_k$ on the doubled bond space
(\leanid{MPSTensor.exists_cutGauge_of_coeff_eq}), and consecutive cuts
satisfy the printed chain relation
\begin{align}
    W_k(C_i\otimes\Id)
     & =(B_i\otimes\Id)W_{k+1}
    \label{eq:pgvwc07_ti_uniqueness_scope_chain_relation}
\end{align}
of lines~1084--1086 (\leanid{MPSTensor.cutGauge_chain}). This is the square
case of Theorem~2, since the cut ranks equal $D^2$ on the window. From the
$D^4+1$ cuts after $L_0,\ldots,L_0+D^4$ sites, the chain-combination and
Kronecker intertwiner results before the proof of Theorem~7 produce the
nonzero intertwiner $R$ with $RC_i=x^{-1}B_iR$
(\leanid{MPSTensor.exists_intertwiner_of_fixedLength_mpv_eq}), with the
intertwiner count corrected as in
\path{docs/paper-gaps/pgvwc07_intertwiner_chain_off_by_one.tex}.

Consequently the printed hypothesis that the open-boundary canonical
representation is unique is not consumed by the formal proof, and it does
not appear in the formal statement. The remaining printed steps (lines
1097--1108) derive $|x|=1$ from the dual fixed point $\Lambda$ and
$RR^\dagger=\Id$ from the single-block property. The formal proof derives
both from condition
(\ref{eq:pgvwc07_ti_uniqueness_scope_finite_injectivity}): block injectivity
makes $R$ invertible and both transfer maps irreducible, and the unitary
gauge theorem for left-canonical irreducible tensors, applied to the
conjugate-transposed families, gives the unitary $U$ and $|x|=1$
(\leanid{MPSTensor.exists_unitaryConj_of_intertwines_of_isNBlkInjective}).
The dual fixed point and the scalar fixed-point condition remain hypotheses
of the formal statement, so no hypothesis of the printed theorem is
strengthened. The root-of-unity condition
(\ref{eq:pgvwc07_ti_uniqueness_scope_root_phase}) follows from the equality
of the two nonzero length-$N$ states.

The printed theorem itself, with its exact conjugacy conclusion, remains
false by the counterexample of
\path{docs/paper-gaps/pgvwc07_quadratic_reconstruction_phase.tex} and
carries no formal declaration; the blueprint states it beside the corrected
theorem.

\section{Tracked consequences}

Blueprint statements that invoke Chapter~9 block separation must state the
strict weight-modulus hypothesis whenever the corresponding \Lean{}
declaration requires it. Downstream CPSV16 uses must not cite these statements
as PGVWC07 canonical-form uniqueness unless a separate formal theorem
establishes finite injectivity, OBC uniqueness, the finite-size bound
(\ref{eq:pgvwc07_ti_uniqueness_scope_finite_size_bound}), and the corrected
unitary-phase conclusion
(\ref{eq:pgvwc07_ti_uniqueness_scope_corrected_unitary_intertwiner}).

\bibliographystyle{alpha}
\bibliography{references}

\end{document}
